\textbf{Seats}
\begin{enumerate}
    \item 8 mins: Intro slides \par
    \textbf{Motivation}
    \item 13 mins: Go back and do exercise 2.2.2 \par
    \textbf{Check piazza}
    \item 20 mins: Finish Example 2.3.1 \par
    \textbf{Check piazza}
    \item 36 mins: Finish Example 2.3.2 \par
    \textbf{Check piazza}
    \item 32 mins: Finish Example 2.3.3 \par
    \textbf{Check piazza}
    \item 40 mins: Finish Example 2.3.4 \par
    \textbf{Check piazza}
    \item 46 mins: Finish Example 2.3.5 (or leave until next time) \par
    \textbf{Check piazza}
    \item 50 mins: Wrap up \par
\end{enumerate}
\textbf{Seats}

\vspace{10mm}
\textbf{Example 2.2.2}
First, find the stationary points $p'(x) = 0$.
\begin{align*}
    p(x) &= 2x^3-3x^2-36x+13\\
    p'(x) &= 6x^2-6x-36\\
    &= 6(x-3)(x+2).
\end{align*}
So, stationary points at $x = \{-2,3\}$.

The function is continuous so we just need to check
\begin{itemize}
    \item The left-hand limit: $x=-4 \quad \rightarrow \quad p(-4) = 2(-64) - 3(16) - 36(-4) + 13 = -19$
    \item The right-hand limit: $x=6 \quad \rightarrow \quad p(6) = 2(216)-3(36)-36(6)+13 = 121$
    \item The first stationary point $x=-2 \quad \rightarrow \quad p(-2) = 2(-8) - 3(4) - 36(-2) + 13 = 57$
    \item The second stationary point $x=3 \quad \rightarrow \quad p(3) = (27)-3(9)-36(3)+13=-68$
\end{itemize}
So, 
\begin{itemize}
    \item Minimum is $p=-68$ at $x=3$.
    \item Maximum is $p = 121$ at $x=6$.
\end{itemize}
Therefore, the minimum of $p(x)$ over the interval $-4 \leq x \leq 6$ occurs at the stationary point $x=3$, where the function takes the value $p(3) = -68$. The maximum occurs at the upper limit of the interval $x=6$, where $p(6) = 121$.

\vspace{10mm}
\textbf{Example 2.3.1}
(i) Expanding the brackets gives 
\begin{align*}
    p(t) &= (t^2+1)(t^2+1)(t^2+1)\\
    &= t^6+3t^4+3t^4+1.
\end{align*} 
Differentiating this polynomial with respect to $t$ gives 
\begin{align*}
    p'(t) = 6t^5 + 12t^3 + 6t,
\end{align*}
which we can simplify to
\begin{align*}
p'(t) &= 6t(t^4+2t^2+1)\\
&=6t(t^2+1)^2
\end{align*}
through factorisation.\\

(ii) Define the inner function as $u=g(t)=t^2+1$, and so the outer function is $p=h(u)=u^3$. Differentiating these functions we get 
\begin{align*}
    h'(u)=3u^2 \qquad \mathrm{and} \qquad g'(t)=2t.
\end{align*} 
Using the chain rule and substituting in these expressions gives
\begin{align*}
    p'(t) &= 3u^2 \times 2t \\
    &= 6t(t^2+1)^2,
\end{align*}
where in the final step we substituted in the expression for $u$.


\vspace{10mm}
\textbf{Example 2.3.2}

Note that $5x-\pi$ is a linear function, and so we can use the rule introduced above. The derivative of $h(u) = \cos u$ is $h'(u) = -\sin u$, and therefore the solution is
\begin{align*}
    r'(x) = -5\sin(5x-\pi).
\end{align*}

\vspace{10mm}
\textbf{Example 2.3.3}

Let's define $u(x) = 3x^2$ and $v(x) = \ln x$ (remember, the order is unimportant). Therefore, $u'(x) = 6x$ and $v'(x) = 1/x$. Using the product rule we get
\begin{align*}
    g'(x) &= u(x)v'(x)+u'(x)v(x) \\
    &= 3x^2 / x + 6x \ln x \\
    &= 3x (1 + 2\ln x).
\end{align*}

\vspace{10mm}
\textbf{Example 2.3.4}

\begin{align*}
    g'(x) &= u(x)v(x)w'(x) + u(x)v'(x)w(x) + u'(x)v(x)w(x) \\
    &= x^2 e^x \cos x + x^2 e^x \sin x + 2x e^x \sin x \\
    &= xe^x(x \cos x + x \sin x + 2\sin x).
\end{align*}

\vspace{10mm}
\textbf{Example 2.3.5}

\begin{align*}
    g(x) &= \frac{\ln(x)}{x^2}\\
    g'(x) &=\frac{u'(x)v(x)-u(x)v'(x)}{v(x)^2}\\
    g'(x) &= \frac{x^2/x-2x\ln(x)}{x^4}\\
    g'(x) &= \frac{1-2\ln(x)}{x^3}
\end{align*}